% GENERATED FILE. Edit canonical lessons, then run npm run generate:books.
% canonical-source-hash: fnv1a64:574c6a3976b2faac
% canonical-lessons: JA-C385-kaigai, JA-C385-gaikoku, JA-C385-kokusai, JA-C385-jisa, JA-C385-mezasu

\chapter{Going Abroad}
\label{ch:ja-going-abroad}
\hlchaptermodality{385}

\begin{chapteropening}
\textbf{By the end of this chapter:} \emph{I can talk about travelling to other countries.}

Complete the last lesson of chapter 385: I can talk about travelling to other countries.
\end{chapteropening}

\section[\texorpdfstring{\ja{かいがい} (kaigai)}{かいがい (kaigai)}]{\ja{かいがい} (kaigai) — overseas; abroad}

\label{lesson:JA-C385-kaigai}



\begin{quote}
Before the new one: say the Japanese for landing (of a plane), then the Japanese for to set off on a journey.
\end{quote}

\subsection*{You'll want to know: \ja{かいがい}}
\textbf{\ja{かいがい}} — \emph{kaigai} — \textquotedblleft{}overseas; abroad\textquotedblright{}.

\begin{quote}
\textbf{\ja{らいねん} \ja{かいがいに} \ja{いきたいです。}} — \emph{rainen kaigai ni ikitai desu} — \textquotedblleft{}I want to go abroad next year.\textquotedblright{}
\end{quote}

\begin{cousinweb}
\textbf{\ja{かいがい}} (\emph{kaigai}) is Sino-Japanese (kango): two Chinese characters read with their on'yomi, readings borrowed from Chinese. They are \emph{kai}, \textquotedblleft{}sea\textquotedblright{} (Middle Chinese \emph{xojX}) and \emph{gai}, \textquotedblleft{}outside\textquotedblright{} (\emph{ngwajH}) (KANJIDIC2; Middle Chinese in Baxter's transcription).
\end{cousinweb}

\begin{grammarlens}[title={What you've built}]
One more word for getting around.
\end{grammarlens}

\subsection*{Your turn}
\begin{itemize}
\raggedright
  \item \emph{Say it:} \emph{kaigai}
  \item \emph{Say it:} \emph{kaigai}, once more
  \item \emph{Say it:} \emph{rainen kaigai ni ikitai desu} — I want to go abroad next year.
  \item \emph{Recall:} say the Japanese for landing (of a plane), then the Japanese for to set off on a journey, then say \emph{kaigai} again
\end{itemize}

\begin{tcolorbox}[breakable,colback=teal!4,colframe=teal!35!black,title={Before you move on}]
What does \ja{かいがい} mean? (\textquotedblleft{}Overseas; abroad\textquotedblright{}.) Say it once more.
\end{tcolorbox}
\section[\texorpdfstring{\ja{がいこく} (gaikoku)}{がいこく (gaikoku)}]{\ja{がいこく} (gaikoku) — a foreign country}

\label{lesson:JA-C385-gaikoku}



\begin{quote}
Before the new one: say the Japanese for to set off on a journey, then the Japanese for overseas.
\end{quote}

\subsection*{You'll want to know: \ja{がいこく}}
\textbf{\ja{がいこく}} — \emph{gaikoku} — \textquotedblleft{}a foreign country\textquotedblright{}.

\begin{quote}
\textbf{\ja{がいこくの} \ja{おかねに} \ja{りょうがえしました。}} — \emph{gaikoku no okane ni ryōgae shimashita} — \textquotedblleft{}I changed my money into foreign currency.\textquotedblright{}
\end{quote}

\begin{cousinweb}
\textbf{\ja{がいこく}} (\emph{gaikoku}) is Sino-Japanese (kango): two Chinese characters read with their on'yomi, readings borrowed from Chinese. They are \emph{gai}, \textquotedblleft{}outside\textquotedblright{} (Middle Chinese \emph{ngwajH}) and \emph{koku}, \textquotedblleft{}country\textquotedblright{} (\emph{kwok}) (KANJIDIC2; Middle Chinese in Baxter's transcription).
\end{cousinweb}

\begin{grammarlens}[title={What you've built}]
One more word for getting around.
\end{grammarlens}

\subsection*{Your turn}
\begin{itemize}
\raggedright
  \item \emph{Say it:} \emph{gaikoku}
  \item \emph{Say it:} \emph{gaikoku}, once more
  \item \emph{Say it:} \emph{gaikoku no okane ni ryōgae shimashita} — I changed my money into foreign currency.
  \item \emph{Recall:} say the Japanese for to set off on a journey, then the Japanese for overseas, then say \emph{gaikoku} again
\end{itemize}

\begin{tcolorbox}[breakable,colback=teal!4,colframe=teal!35!black,title={Before you move on}]
What does \ja{がいこく} mean? (\textquotedblleft{}A foreign country\textquotedblright{}.) Say it once more.
\end{tcolorbox}
\section[\texorpdfstring{\ja{こくさい} (kokusai)}{こくさい (kokusai)}]{\ja{こくさい} (kokusai) — international}

\label{lesson:JA-C385-kokusai}



\begin{quote}
Before the new one: say the Japanese for overseas, then the Japanese for a foreign country.
\end{quote}

\subsection*{You'll want to know: \ja{こくさい}}
\textbf{\ja{こくさい}} — \emph{kokusai} — \textquotedblleft{}international\textquotedblright{}.

\begin{quote}
\textbf{\ja{この} \ja{くうこうは} \ja{こくさい} \ja{くうこうです。}} — \emph{kono kūkō wa kokusai kūkō desu} — \textquotedblleft{}This airport is an international airport.\textquotedblright{}
\end{quote}

\begin{cousinweb}
\textbf{\ja{こくさい}} (\emph{kokusai}) is Sino-Japanese (kango): two Chinese characters read with their on'yomi, readings borrowed from Chinese. They are \emph{koku}, \textquotedblleft{}country\textquotedblright{} (Middle Chinese \emph{kwok}) and \emph{sai}, \textquotedblleft{}edge\textquotedblright{} (\emph{tsjejH}) (KANJIDIC2; Middle Chinese in Baxter's transcription).
\end{cousinweb}

\begin{grammarlens}[title={What you've built}]
One more word for getting around.
\end{grammarlens}

\subsection*{Your turn}
\begin{itemize}
\raggedright
  \item \emph{Say it:} \emph{kokusai}
  \item \emph{Say it:} \emph{kokusai}, once more
  \item \emph{Say it:} \emph{kono kūkō wa kokusai kūkō desu} — This airport is an international airport.
  \item \emph{Recall:} say the Japanese for overseas, then the Japanese for a foreign country, then say \emph{kokusai} again
\end{itemize}

\begin{tcolorbox}[breakable,colback=teal!4,colframe=teal!35!black,title={Before you move on}]
What does \ja{こくさい} mean? (\textquotedblleft{}International\textquotedblright{}.) Say it once more.
\end{tcolorbox}
\section[\texorpdfstring{\ja{じさ} (jisa)}{じさ (jisa)}]{\ja{じさ} (jisa) — a time difference}

\label{lesson:JA-C385-jisa}



\begin{quote}
Before the new one: say the Japanese for a foreign country, then the Japanese for international.
\end{quote}

\subsection*{You'll want to know: \ja{じさ}}
\textbf{\ja{じさ}} — \emph{jisa} — \textquotedblleft{}a time difference\textquotedblright{}.

\begin{quote}
\textbf{\ja{じさで} \ja{つかれました。}} — \emph{jisa de tsukaremashita} — \textquotedblleft{}The time difference left me tired.\textquotedblright{}
\end{quote}

\begin{cousinweb}
\textbf{\ja{じさ}} (\emph{jisa}) is Sino-Japanese (kango): two Chinese characters read with their on'yomi, readings borrowed from Chinese. They are \emph{ji}, \textquotedblleft{}time\textquotedblright{} (Middle Chinese \emph{dzyi}) and \emph{sa}, \textquotedblleft{}difference\textquotedblright{} (\emph{tsrhae}) (KANJIDIC2; Middle Chinese in Baxter's transcription).
\end{cousinweb}

\begin{grammarlens}[title={What you've built}]
One more word for getting around.
\end{grammarlens}

\subsection*{Your turn}
\begin{itemize}
\raggedright
  \item \emph{Say it:} \emph{jisa}
  \item \emph{Say it:} \emph{jisa}, once more
  \item \emph{Say it:} \emph{jisa de tsukaremashita} — The time difference left me tired.
  \item \emph{Recall:} say the Japanese for a foreign country, then the Japanese for international, then say \emph{jisa} again
\end{itemize}

\begin{tcolorbox}[breakable,colback=teal!4,colframe=teal!35!black,title={Before you move on}]
What does \ja{じさ} mean? (\textquotedblleft{}A time difference\textquotedblright{}.) Say it once more.
\end{tcolorbox}
\section[\texorpdfstring{\ja{めざす} (mezasu)}{めざす (mezasu)}]{\ja{めざす} (mezasu) — to aim for; to make for}

\label{lesson:JA-C385-mezasu}



\begin{quote}
Before the new one: say the Japanese for international, then the Japanese for a time difference.
\end{quote}

\subsection*{You'll want to know: \ja{めざす}}
\textbf{\ja{めざす}} — \emph{mezasu} — \textquotedblleft{}to aim for; to make for\textquotedblright{}.

\begin{quote}
\textbf{\ja{ふねは} \ja{みなみの} \ja{しまを} \ja{めざしています。}} — \emph{fune wa minami no shima o mezashite imasu} — \textquotedblleft{}The ship is making for the southern islands.\textquotedblright{}
\end{quote}

Polite \textbf{\ja{めざします}}, \emph{mezashimasu}; a godan verb. Past (polite): \textbf{\ja{ふねが} \ja{しまを} \ja{めざしました。}}, \emph{fune ga shima o mezashimashita}.

\begin{cousinweb}
\textbf{\ja{めざす}} (\emph{mezasu}) is a native Japanese word (wago), inherited rather than borrowed from Chinese, built from \emph{me}, \textquotedblleft{}eye\textquotedblright{} and \emph{zasu}, \textquotedblleft{}pointing\textquotedblright{} (from \emph{sasu}); its kanji were attached later as kun'yomi (JMdict).
\end{cousinweb}

\begin{grammarlens}[title={What you've built}]
One more word for getting around.
\end{grammarlens}

\subsection*{Your turn}
\begin{itemize}
\raggedright
  \item \emph{Say it:} \emph{mezasu}
  \item \emph{Say it:} \emph{mezasu}, once more
  \item \emph{Say it:} \emph{fune wa minami no shima o mezashite imasu} — The ship is making for the southern islands.
  \item \emph{Recall:} say the Japanese for international, then the Japanese for a time difference, then say \emph{mezasu} again
\end{itemize}

\begin{tcolorbox}[breakable,colback=teal!4,colframe=teal!35!black,title={Before you move on}]
What does \ja{めざす} mean? (\textquotedblleft{}To aim for; to make for\textquotedblright{}.) Say it once more.
\end{tcolorbox}
